\chapter{Exponencial e logaritmo}\label{ch:g12:exp}

A \omterm{thm:g12:exp:existence}{função exponencial} é a única \omterm{def:g11:func:function}{função} igual à própria \omterm{def:g12:deriv:derivative}{derivada} e que vale
$1$ em $0$. Ela converte somas em produtos; sua inversa, o \omterm{def:g12:exp:ln}{logaritmo
natural}, converte produtos em somas. Juntas, descrevem todo fenômeno cuja
taxa de variação é proporcional ao seu tamanho: decaimento radioativo,
crescimento populacional, juros compostos.

\section{A função exponencial}

\begin{theorem}[Existência e unicidade]\label{thm:g12:exp:existence}
Existe uma única \omterm{def:g11:func:function}{função} \omterm{def:g12:deriv:derivative}{derivável} $f \colon \R \to \R$ tal que
\[
f' = f \qquad\text{e}\qquad f(0) = 1 .
\]
Ela é chamada de \emph{função exponencial}\index{exponencial} e escrita
$\exp$, ou $x \mapsto \eu^x$.
\end{theorem}

\begin{proof}[Demonstração da unicidade]
Primeiro, uma \omterm{def:g11:func:function}{função} assim nunca se anula. De fato, seja
$g(x) = f(x)f(-x)$; então
\[
g'(x) = f'(x)f(-x) - f(x)f'(-x) = f(x)f(-x) - f(x)f(-x) = 0,
\]
de modo que $g$ é constante igual a $g(0) = 1$: para todo $x$,
$f(x)f(-x) = 1$ e, em particular, $f(x) \neq 0$.

Agora sejam $f_1, f_2$ duas soluções e $h = \dfrac{f_1}{f_2}$ (legítimo,
já que $f_2$ nunca se anula). Então
\[
h' = \frac{f_1' f_2 - f_1 f_2'}{f_2^2} = \frac{f_1 f_2 - f_1 f_2}{f_2^2} = 0,
\]
de modo que $h$ é constante igual a $h(0) = 1$, \emph{isto é},
$f_1 = f_2$.

A \emph{existência} é admitida neste nível (pode ser obtida pela \omterm{def:g12:seq:sequence}{sequência}
do \cref{exo:g12:seq:10}, ou como inversa do \omterm{def:g12:exp:ln}{logaritmo} construído por
integração no \cref{ch:g12:integ}).
\end{proof}

\begin{proposition}[Equação funcional]\label{prop:g12:exp:funceq}
Para todos $x, y \in \R$ e $n \in \Z$:
\[
\eu^{x+y} = \eu^x\,\eu^y, \qquad
\eu^{-x} = \frac{1}{\eu^x}, \qquad
\eu^{x-y} = \frac{\eu^x}{\eu^y}, \qquad
\bigl(\eu^{x}\bigr)^n = \eu^{nx}.
\]
Além disso, $\eu^x > 0$ para todo $x$.
\end{proposition}

\begin{proof}
Fixe $y$ e considere $\varphi(x) = \dfrac{\exp(x+y)}{\exp(x)\exp(y)}$. Seu
numerador e seu denominador, como funções de $x$, são ambos soluções de
$f' = f$ a menos da constante $\exp(y)$; derivar $\varphi$ diretamente
(regra do quociente) dá $\varphi' = 0$, de modo que
$\varphi \equiv \varphi(0) = 1$, o que demonstra a primeira identidade.
Tomar $y = -x$ dá a segunda (com o valor $\eu^0 = 1$), e a terceira segue.
A quarta é uma indução a partir da primeira para $n \geq 0$, estendida a
$n < 0$ pela segunda.

Positividade: $\eu^x = \left(\eu^{x/2}\right)^2 \geq 0$ e $\eu^x \neq 0$
(mostrado no \cref{thm:g12:exp:existence}), logo $\eu^x > 0$.
\end{proof}

\begin{proposition}[Variação e limites]\label{prop:g12:exp:variations}
A exponencial é estritamente \omterm{def:g12:seq:monotonic}{crescente}, \omterm{def:g12:deriv:convex}{convexa}, e
\[
\lim_{x\to-\infty} \eu^x = 0, \qquad
\lim_{x\to+\infty} \eu^x = +\infty .
\]
\end{proposition}

\begin{proof}
$(\exp)' = \exp > 0$ dá o crescimento estrito; $(\exp)'' = \exp > 0$ dá a
\omterm{def:g12:deriv:convex}{convexidade}. Pela \omterm{def:g12:deriv:convex}{convexidade}, $\eu^x \geq 1 + x$ (\omterm{def:g12:deriv:tangent}{tangente} em $0$, veja o
\cref{ex:g12:deriv:ineq}), de modo que $\eu^x \to +\infty$ quando
$x \to +\infty$, por comparação. Então
$\eu^{x} = 1/\eu^{-x} \to 0$ quando $x \to -\infty$.
\end{proof}

\begin{omfigure}
\begin{tikzpicture}
\begin{axis}[omaxis,
  width=8.8cm, height=6cm,
  xmin=-3.3, xmax=2.4, ymin=-2.4, ymax=7.8,
  xtick={-2,-1,1,2}, ytick={1}]
  \addplot[omProp, thick, domain=-3.2:2.3] {1 + x};
  \addplot[omDef, thick, domain=-3.2:2.0] {exp(x)};
  \fill (axis cs:0,1) circle (1.5pt);
  \node[omDef, font=\small, anchor=east] at (axis cs:1.35,5.2) {$\eu^x$};
  \node[omProp, font=\small, anchor=north west] at (axis cs:1.05,1.6) {$y = 1+x$};
\end{axis}
\end{tikzpicture}

{\small A exponencial: estritamente \omterm{def:g12:seq:monotonic}{crescente}, \omterm{def:g12:deriv:convex}{convexa}, com
$\lim_{-\infty} \eu^x = 0$. Ela fica acima de sua \omterm{def:g12:deriv:tangent}{tangente} em $0$:
$\eu^x \geq 1 + x$.}
\end{omfigure}

\begin{theorem}[Comparação de crescimentos]\label{thm:g12:exp:growth}
\index{comparação de crescimentos}
Para todo inteiro $n \geq 1$,
\[
\lim_{x\to+\infty} \frac{\eu^x}{x^n} = +\infty,
\qquad
\lim_{x\to-\infty} x^n \eu^x = 0 .
\]
Em palavras: a exponencial vence toda potência de $x$.
\end{theorem}

\begin{proof}
Para $n = 1$: aplicando $\eu^t \geq 1 + t \geq t$ em $t = x/2$,
\[
\frac{\eu^x}{x} = \frac{\left(\eu^{x/2}\right)^2}{x}
\geq \frac{(x/2)^2}{x} = \frac{x}{4} \xrightarrow[x\to+\infty]{} +\infty .
\]
Para $n$ geral, escreva
\[
\frac{\eu^x}{x^n}
= \frac{1}{n^n}\left(\frac{\eu^{x/n}}{x/n}\right)^{n}
\xrightarrow[x\to+\infty]{} +\infty
\]
pelo caso $n=1$ e por composição. O limite em $-\infty$ segue da
substituição $x \mapsto -x$:
$\abs{x^n \eu^x} = \frac{\abs{x}^n}{\eu^{\abs x}} \to 0$.
\end{proof}

\begin{omfigure}
\begin{tikzpicture}
\begin{axis}[omaxis,
  width=8.8cm, height=6cm,
  xmin=-0.4, xmax=5.6, ymin=-12, ymax=155,
  xtick={1,2,3,4,5}, ytick={50,100,150}]
  \addplot[omThm, thick, domain=0:5.4] {x^2};
  \addplot[omProp, thick, domain=0:5.3] {x^3};
  \addplot[omDef, thick, domain=-0.4:5.0] {exp(x)};
  \fill (axis cs:4.536,93.4) circle (1.5pt);
  \draw[black, dashed, thin] (axis cs:4.536,0) -- (axis cs:4.536,93.4);
\end{axis}
\end{tikzpicture}

{\small A exponencial (azul) acaba vencendo toda potência: ela ultrapassa
$x^2$ (vermelho) logo no início e $x^3$ (laranja) em $x \approx 4.54$.}
\end{omfigure}

\section{O logaritmo natural}

\begin{definition}[Logaritmo natural]\label{def:g12:exp:ln}
A exponencial é \omterm{def:g12:limcont:continuity}{contínua} e estritamente \omterm{def:g12:seq:monotonic}{crescente} de $\R$ sobre
$\intoo{0}{+\infty}$; pelo teorema da bijeção
(\cref{thm:g12:limcont:bijection}), para todo $y > 0$ a \omterm{def:g10:algebra:equation}{equação}
$\eu^x = y$ tem solução única. Essa solução é o \emph{logaritmo
natural}\index{logaritmo} de $y$, escrito $\ln y$. Assim,
\[
\text{para } x \in \R,\ y > 0: \qquad y = \eu^x \iff x = \ln y .
\]
Em particular, $\ln 1 = 0$, $\ln \eu = 1$, e $\eu^{\ln y} = y$,
$\ln(\eu^x) = x$.
\end{definition}

\begin{omfigure}
\begin{tikzpicture}
\begin{axis}[omaxis,
  width=8.6cm, axis equal image,
  xmin=-4.3, xmax=6.3, ymin=-4.3, ymax=6.3,
  xtick={1}, ytick={1}]
  \addplot[black, dashed, domain=-3.9:5.9] {x};
  \addplot[omDef, thick, domain=-4.1:1.79] {exp(x)};
  \addplot[omThm, thick, domain=0.0167:5.9, samples=120] {ln(x)};
  \node[omDef, font=\small, anchor=east] at (axis cs:0.95,2.6) {$\eu^x$};
  \node[omThm, font=\small, anchor=north] at (axis cs:4.3,1.25) {$\ln x$};
  \node[font=\small, anchor=west, rotate=45] at (axis cs:4.6,4.85) {$y=x$};
\end{axis}
\end{tikzpicture}

{\small As curvas de $\exp$ e de $\ln$ são espelhos uma da outra na reta
$y = x$: o ponto $(x, \eu^x)$ reflete-se em $(\eu^x, x)$.}
\end{omfigure}

\begin{proposition}[Propriedades algébricas]\label{prop:g12:exp:lnalg}
Para todos $a, b > 0$ e $n \in \Z$:
\[
\ln(ab) = \ln a + \ln b, \quad
\ln\frac{1}{a} = -\ln a, \quad
\ln\frac{a}{b} = \ln a - \ln b, \quad
\ln(a^n) = n \ln a, \quad
\ln\sqrt{a} = \tfrac12 \ln a .
\]
\end{proposition}

\begin{proof}
$\eu^{\ln a + \ln b} = \eu^{\ln a}\eu^{\ln b} = ab$, e tomar $\ln$ dos dois
lados dá a primeira identidade. As demais seguem pelo mesmo mecanismo, a
partir das identidades correspondentes da \cref{prop:g12:exp:funceq}.
\end{proof}

\begin{proposition}[Propriedades analíticas]\label{prop:g12:exp:lnanalytic}
A \omterm{def:g11:func:function}{função} $\ln$ é \omterm{def:g12:deriv:derivative}{derivável} em $\intoo{0}{+\infty}$, com
\[
(\ln x)' = \frac{1}{x},
\]
estritamente \omterm{def:g12:seq:monotonic}{crescente}, \omterm{def:g12:deriv:convex}{côncava}, e
$\lim\limits_{x\to0^+} \ln x = -\infty$,
$\lim\limits_{x\to+\infty} \ln x = +\infty$. Além disso, para todo inteiro
$n \geq 1$,
\[
\lim_{x\to+\infty} \frac{\ln x}{x^{1/n}} = 0, \qquad
\lim_{x\to0^+} x \ln x = 0 .
\]
\end{proposition}

\begin{proof}
A derivabilidade da \omterm{def:g11:func:function}{função} inversa é admitida neste nível; admitindo-a,
derive a identidade $\eu^{\ln x} = x$ pela regra da cadeia:
$(\ln x)'\,\eu^{\ln x} = 1$, logo $(\ln x)' = \frac{1}{x} > 0$, donde o
crescimento estrito, e $(\ln x)'' = -\frac1{x^2} < 0$, donde a
concavidade. Os limites em $0^+$ e em $+\infty$ espelham os de $\exp$ pela
bijeção. Para a \omterm{thm:g12:exp:growth}{comparação de crescimentos}, substitua $x = \eu^{t}$:
$\frac{\ln x}{x} = \frac{t}{\eu^t} \to 0$ quando $t \to +\infty$, pelo
\cref{thm:g12:exp:growth}, e do mesmo modo para os demais limites com
$x = \eu^{-t}$, o que dá $x\ln x = -t\eu^{-t} \to 0$.
\end{proof}

\begin{method}[Resolver equações com $\exp$ e $\ln$]\label{met:g12:exp:equations}
As duas funções são estritamente \omterm{def:g12:seq:monotonic}{crescentes}, de modo que podem ser
aplicadas aos (ou retiradas dos) dois lados de uma \omterm{def:g10:algebra:equation}{equação} ou inequação
sem mudar seu sentido:
\[
\eu^{u} = \eu^{v} \iff u = v, \qquad
\eu^{u} \leq \eu^{v} \iff u \leq v,
\]
e do mesmo modo para $\ln$ em quantidades positivas. \emph{Verifique
sempre os \omterm{def:g11:func:function}{domínios} primeiro} ($\ln$ exige argumentos positivos). Para
\omterm{def:g10:algebra:equation}{equações} como $a^x = b$ com $a>0$, $a \neq 1$, reescreva
$a^x = \eu^{x\ln a}$ e resolva $x = \frac{\ln b}{\ln a}$.
\end{method}

\begin{example}\label{ex:g12:exp:halflife}
Uma \omterm{def:g10:proba:sample}{amostra} radioativa decai segundo $N(t) = N_0\,\eu^{-\lambda t}$. Sua
\emph{meia-vida} $T$ satisfaz $N(T) = N_0/2$, \emph{isto é},
$\eu^{-\lambda T} = \frac12$, logo $T = \frac{\ln 2}{\lambda}$: a
meia-vida não depende da quantidade inicial.
\end{example}

\section{Exercícios}

\begin{exercise}[$\star$]\label{exo:g12:exp:1}
Simplifique $\dfrac{\eu^{3x}\,\eu^{-x+1}}{\eu^{x}}$ e
$\ln\!\left(\dfrac{\eu^2\sqrt{\eu}}{\eu^{-3}}\right)$.
\end{exercise}

\begin{exercise}[$\star$]\label{exo:g12:exp:2}
Resolva em $\R$:
\[
\text{(a) } \eu^{2x} - 3\eu^{x} + 2 = 0;
\qquad
\text{(b) } \ln(x - 1) + \ln(x + 2) = \ln 4 .
\]
\end{exercise}

\begin{exercise}[$\star$]\label{exo:g12:exp:3}
Calcule os limites:
\[
\lim_{x\to+\infty} \frac{\eu^x - x^2}{\eu^x + 1}, \qquad
\lim_{x\to+\infty} \frac{\ln x}{\sqrt x}, \qquad
\lim_{x\to0^+} x^2 \ln x, \qquad
\lim_{x\to+\infty} \bigl(x - \ln x\bigr).
\]
\end{exercise}

\begin{exercise}[$\star\star$]\label{exo:g12:exp:4}
Estude a \omterm{def:g11:func:function}{função} $f(x) = x\,\eu^{-x}$ em $\R$: variação, limites, \omterm{def:g12:deriv:extremum}{extremo};
mostre que sua curva tem um \omterm{def:g12:deriv:inflection}{ponto de inflexão} e dê suas \omterm{def:g10:coordgeom:system}{coordenadas}.
\end{exercise}

\begin{exercise}[$\star\star$]\label{exo:g12:exp:5}
Mostre que, para todo $x > -1$, $\ln(1 + x) \leq x$, com igualdade somente
em $x = 0$. Deduza que, para todo $n \geq 1$,
\[
\left(1 + \frac{1}{n}\right)^n \leq \eu \leq \left(1 - \frac{1}{n+1}\right)^{-(n+1)} .
\]
\end{exercise}

\begin{exercise}[$\star\star$]\label{exo:g12:exp:6}
Um capital $C_0$ é aplicado a uma taxa anual de $3\%$, com juros
capitalizados a cada ano.
\begin{enumerate}
  \item Exprima o capital $C_n$ após $n$ anos.
  \item Depois de quantos anos o capital dobra? Dê a resposta exata usando
        $\ln$ e depois um valor numérico.
  \item Compare com a \omterm{def:g10:numbers:approx}{aproximação} ``$70$ dividido pela taxa em
        porcentagem'', usada pelos banqueiros.
\end{enumerate}
\end{exercise}

\begin{exercise}[$\star\star$]\label{exo:g12:exp:7}
Resolva a inequação $\eu^{2x} - \eu^{x+1} > 0$ e depois a inequação
$\ln(x^2 - 1) \leq \ln(x + 5)$.
\end{exercise}

\begin{exercise}[$\star\star\star$]\label{exo:g12:exp:8}
Seja $f(x) = \dfrac{\eu^x}{x}$ para $x > 0$.
\begin{enumerate}
  \item Estude a variação de $f$ em $\intoo{0}{+\infty}$ e dê seu \omterm{def:g10:functions:extrema}{mínimo}.
  \item Para quais valores de $k$ a \omterm{def:g10:algebra:equation}{equação} $\eu^x = kx$ tem $0$, $1$ ou
        $2$ soluções em $\intoo{0}{+\infty}$?
\end{enumerate}
\end{exercise}

\begin{exercise}[$\star\star\star$]\label{exo:g12:exp:9}
Para $n \geq 1$, seja $u_n = \left(1 + \frac1n\right)^n$.
\begin{enumerate}
  \item Usando o \cref{exo:g12:exp:5}, mostre que $(u_n)$ é \omterm{def:g12:seq:bounded}{limitada
        superiormente} por~$\eu$.
  \item Mostre que $\ln u_n = n \ln\left(1 + \frac1n\right) \to 1$ e deduza
        que $u_n \to \eu$. (Sugestão: reconheça um quociente de diferenças
        de $\ln$ em $1$.)
\end{enumerate}
\end{exercise}

\section{Problema: O logaritmo domestica o mundo}

\begin{problem}[{Problema de fim de semana --- tempos de duplicação e
a regra de 72, as escalas de terremotos, ácidos e pianos, e a
constante $\eu$ deixando digitais em toda parte}]
\label{pb:g12:exp:1}
Tudo o que cresce por uma \emph{porcentagem} fixa cresce
exponencialmente --- poupanças, bactérias, epidemias --- e tudo o
que abrange potências de dez demais para se apreender --- energias
de terremotos, acidez, intensidades sonoras --- é domesticado por
um \omterm{def:g12:exp:ln}{logaritmo}. Este problema calcula tempos de duplicação e desmonta
a regra de 72 dos banqueiros, lê as escalas logarítmicas do mundo e
recolhe as digitais que o número $\eu$ deixa em cada cena de crime
(\cref{prop:g12:exp:funceq}, \cref{thm:g12:exp:growth},
\cref{met:g12:exp:equations}).

\medskip
\noindent\textbf{Parte I --- Fluência.}
\begin{enumerate}
  \item Resolva: $\eu^{2x} = 5$; \quad $\ln(3x - 1) = 2$; \quad
        $\eu^{2x} - 3\eu^x + 2 = 0$ (um \omterm{def:g11:quad:quadratic}{trinômio do segundo grau}
        disfarçado).
  \item Simplifique: $\dfrac{\ln 8}{\ln 2}$; \quad
        $\ln\!\left(\eu^3 \sqrt{\eu}\right)$; \quad
        $\eu^{\ln 5 - \ln 2}$.
  \item Demonstre a desigualdade-mãe: $\eu^x \geq 1 + x$ para
        todo real $x$ (estude $f(x) = \eu^x - 1 - x$) e deduza
        $\ln(1 + u) \leq u$ para $u > -1$.
  \item Calcule $\lim_{x \to +\infty} x^2 \eu^{-x}$,
        $\lim_{x \to +\infty} \frac{\ln x}{x}$
        (\cref{thm:g12:exp:growth}) e
        $\lim_{x \to 0} \frac{\eu^x - 1}{x}$ (reconheça uma
        \omterm{def:g12:deriv:derivative}{derivada}).
  \item Estude $f(x) = x \ln x$ em $\intoo{0}{+\infty}$:
        variação, \omterm{def:g10:functions:extrema}{mínimo} e o limite em $0^+$ (admitido:
        $x \ln x \to 0$). Essa pequena \omterm{def:g11:func:function}{função} mede informação e
        entropia ao longo dos volumes de graduação.
\end{enumerate}

\medskip
\noindent\textbf{Parte II --- Tempos de duplicação e a regra de
72.}
\begin{enumerate}[resume]
  \item Uma poupança rende $3\,\%$ ao ano. Resolva
        $1.03^n = 2$: quanto tempo para dobrar o capital?
  \item Os banqueiros estimam o tempo de duplicação como
        $\frac{72}{\text{taxa em }\%}$. Teste a regra a
        $3\,\%$, $6\,\%$ e $9\,\%$ contra o valor exato
        $\frac{\ln 2}{\ln(1 + r)}$ e explique-a: para $r$
        pequeno, $\ln(1 + r) \approx r$ (a desigualdade da
        questão 3 é metade da história), de modo que a constante
        exata é $100 \ln 2 \approx 69.3$ --- por que os
        banqueiros preferem $72$?
  \item Uma bactéria se divide a cada $20$ minutos:
        $p(t) = 2^{t/20}$ ($t$ em minutos). Reescreva isso como
        $\eu^{\lambda t}$ e calcule $p$ após $24$ horas. A
        resposta (mais de $10^{21}$) prova o quê sobre o modelo
        --- e o que detém as colônias reais?
  \item A cafeína deixa o organismo com meia-vida de cerca de
        $5$ horas. Após um café de $100$ mg às 15h, quanto resta
        às 23h? A que horas o valor cai abaixo de $10$ mg? (A
        insônia tem um \omterm{def:g12:exp:ln}{logaritmo}.)
  \item Um euro a $100\,\%$ de juros anuais: calcule o capital no
        fim do ano com capitalização anual, mensal e diária, e dê
        o teto que o \cref{exo:g12:exp:9} demonstrou ser
        inquebrável. Que constante famosa é o limite da pura
        ganância?
  \item Por que os cientistas plotam $\ln(\text{casos})$ contra o
        tempo na fase inicial de uma epidemia? O que uma reta
        nesse \omterm{def:g10:functions:graph}{gráfico} revela, e o que seu \omterm{def:g10:reffunc:affine}{coeficiente angular}
        mede?
\end{enumerate}

\medskip
\noindent\textbf{Parte III --- As escalas logarítmicas do
mundo.}
(Escreva $\log_{10} x = \frac{\ln x}{\ln 10}$.)
\begin{enumerate}[resume]
  \item Nível sonoro em decibéis: $L = 10 \log_{10}(I/I_0)$. Uma
        conversa mede $60$ dB e um show de rock, $120$ dB: por
        que fator as intensidades sonoras diferem?
  \item As magnitudes de terremotos sobem $1$ quando a amplitude
        sísmica é multiplicada por $10$, e a energia liberada
        escala como amplitude$^{3/2}$. Compare terremotos de
        magnitude $5$ e $7$: razão das amplitudes e depois razão
        das energias.
  \item Química: $\text{pH} = -\log_{10}[\mathrm{H^+}]$. O suco
        de limão tem pH $2$ e o leite, pH $6.5$: qual é a razão
        entre suas concentrações de ácido?
  \item Música: cada oitava dobra a \omterm{def:g10:stats:series}{frequência}. Do lá mais grave
        do piano ($27.5$ Hz) ao dó mais agudo ($4\,186$ Hz),
        quantas oitavas o teclado abrange? E por que nossos
        sentidos --- audição, visão, percepção de tremores ---
        preferem escalas logarítmicas? (Uma frase.)
  \item A régua de cálculo, a calculadora dos engenheiros por
        $350$ anos: duas réguas graduadas de modo que o
        \emph{comprimento} até a marca $x$ seja proporcional a
        $\ln x$. Explique como deslizar uma régua sobre a outra
        multiplica números e nomeie a identidade da
        \cref{prop:g12:exp:lnalg} que faz o trabalho.
\end{enumerate}

\medskip
\noindent\textbf{Parte IV --- As digitais de $\eu$.}
\begin{enumerate}[resume]
  \item Do \cref{exo:g12:exp:8}: a \omterm{def:g10:algebra:equation}{equação} $\eu^x = kx$
        ($x > 0$) tem $0$, $1$ ou $2$ soluções conforme a posição
        de $k$ em relação a um limiar. Reenuncie o resultado ---
        e verifique o fato geométrico por trás dele: a reta
        $y = \eu x$ é exatamente a \omterm{def:g12:deriv:tangent}{tangente} à exponencial que
        passa pela origem.
  \item A constante do quase: com $n$ bilhetes de loteria de
        \omterm{def:g10:proba:distribution}{probabilidade} de vitória $\frac1n$ cada,
        $\P(\text{nenhuma vitória}) = \left(1 -
        \frac1n\right)^n$. Calcule seu limite via
        $n \ln\!\left(1 - \frac1n\right)$ (use o limite padrão
        $\frac{\ln(1 + u)}{u} \to 1$) e reconcilie com o
        misterioso $0.368$ do \cref{pb:g11:binom:1}.
  \item A regra dos $37\,\%$: para escolher o melhor entre $n$
        candidatos entrevistados em ordem aleatória (sem voltar
        atrás), a estratégia ótima rejeita os primeiros
        $\frac{n}{\eu} \approx 37\,\%$ e depois aceita o
        primeiro candidato melhor que todos os anteriores ---
        tendo sucesso com \omterm{def:g10:proba:distribution}{probabilidade} de cerca de
        $\frac1\eu$. De onde vêm os dois $\frac1\eu$ das questões
        18 e 19 --- enuncie o mecanismo comum (muitos \omterm{def:g11:prob:model}{eventos}
        independentes de pequena chance) --- e calcule
        $\frac1\eu$ com três casas decimais.
  \item Final --- o retrato de $\eu$: o teto da capitalização
        (questão 10); a base cuja \omterm{def:g12:deriv:tangent}{tangente} em $0$ tem \omterm{def:g10:reffunc:affine}{coeficiente
        angular} exatamente $1$ (questões 3 e 4); o crescimento
        que nenhum polinômio alcança (questão 4); a constante dos
        quases e da parada ótima (questões 18--19); e sua
        inversa, o $\ln$, que transforma produtos em somas
        (questão 16) e décadas de grandeza em centímetros (Parte
        III). Uma frase para cada.
\end{enumerate}
\end{problem}
